\chapter{Hardware Reference}

This chapter lists only hardware areas and registers that the project touches directly or through GBDK API calls present in the code.

\section{Memory Map Used by the Project}

\begin{table}[h]
\centering
\small
\begin{tabular}{p{2.6cm}p{1.7cm}p{2.6cm}p{6.4cm}}
\toprule
\textbf{Range} & \textbf{Size} & \textbf{Name} & \textbf{Project role} \\
\midrule
\texttt{\$0000--\$3FFF} & 16 KiB & ROM0 & Fixed ROM area generated by GBDK and \texttt{lcc}. \\
\texttt{\$4000--\$7FFF} & 16 KiB & ROMX & Switchable MBC5 ROM banks; Makefile reserves 8 banks. \\
\texttt{\$8000--\$97FF} & 6 KiB & VRAM tile data & Background, window, and sprite tile patterns loaded by \texttt{set\_bkg\_data} and \texttt{set\_sprite\_data}. \\
\texttt{\$9800--\$9BFF} & 1 KiB & BG map 0 & Visible 20x18 tilemaps written by \texttt{set\_bkg\_tiles}, \texttt{set\_win\_tiles}, and \texttt{set\_bkg\_tile\_xy}. \\
\texttt{\$9C00--\$9FFF} & 1 KiB & BG map 1 & Available through GBDK; no direct project selection is found. \\
\texttt{\$A000--\$BFFF} & 8 KiB & External SRAM & Battery-backed save area. \texttt{save\_data\_t} starts at \texttt{\$A000}. \\
\texttt{\$C000--\$DFFF} & 8 KiB & WRAM & C globals, static module state, stack, and runtime data such as \texttt{g\_lobby}, \texttt{g\_tm}, \texttt{g\_tetris}, and \texttt{tetris\_grid}. \\
\texttt{\$FE00--\$FE9F} & 160 B & OAM & Sprite attributes changed by \texttt{move\_sprite} and \texttt{set\_sprite\_tile}. \\
\texttt{\$FF00} & 1 B & JOYP & Joypad state read through \texttt{joypad()}. \\
\texttt{\$FF04} & 1 B & DIV & RNG seed source through \texttt{DIV\_REG}. \\
\texttt{\$FF10--\$FF14} & 5 B & APU channel 1 & Transition and dialogue sound effects. \\
\texttt{\$FF24--\$FF26} & 3 B & APU control & Master volume, routing, and sound enable. \\
\texttt{\$FF40} & 1 B & LCDC & Display, sprite, background, and window enables through GBDK macros. \\
\texttt{\$FF47} & 1 B & BGP & DMG background palette used by fade transitions. \\
\texttt{\$FF4A--\$FF4B} & 2 B & WY/WX & Window position changed by \texttt{move\_win}. \\
\texttt{\$FFFF} & 1 B & IE & VBlank interrupt enable handled by GBDK after \texttt{add\_VBL} and \texttt{enable\_interrupts}. \\
\bottomrule
\end{tabular}
\caption{Memory ranges used by GameBoyGame}
\end{table}

\section{Hardware Registers}

\begin{table}[h]
\centering
\small
\begin{tabular}{p{1.8cm}p{1.9cm}p{1.3cm}p{8.3cm}}
\toprule
\textbf{Address} & \textbf{Register} & \textbf{Access} & \textbf{Bits and project use} \\
\midrule
\texttt{\$FF00} & JOYP/P1 & Read & \texttt{joypad()} returns \texttt{J\_LEFT}, \texttt{J\_RIGHT}, \texttt{J\_UP}, \texttt{J\_DOWN}, \texttt{J\_A}, \texttt{J\_B}, \texttt{J\_START}, and \texttt{J\_SELECT}. \\
\texttt{\$FF04} & DIV & Read & \texttt{random\_init} multiplies \texttt{DIV\_REG} by the first nonzero key mask to seed \texttt{initrand}. \\
\texttt{\$FF10} & NR10 & Write & \texttt{0x00}; channel 1 sweep disabled. \\
\texttt{\$FF11} & NR11 & Write & \texttt{0x81}; channel 1 duty/length value for short effects. \\
\texttt{\$FF12} & NR12 & Write & \texttt{0x43}; initial volume 4, decreasing envelope, period 3. \\
\texttt{\$FF13} & NR13 & Write & \texttt{0x90} for dialogue typing, or \texttt{0x30}, \texttt{0x50}, \texttt{0x70}, \texttt{0x90} for transition tones. \\
\texttt{\$FF14} & NR14 & Write & \texttt{0x86}; trigger bit set and frequency high bits set to 6. \\
\texttt{\$FF24} & NR50 & Write & \texttt{0xFF}; max left/right output volume and VIN bits enabled. \\
\texttt{\$FF25} & NR51 & Write & \texttt{0xFF}; all channels routed to both outputs. \\
\texttt{\$FF26} & NR52 & Write & \texttt{0x80}; sound controller enabled. \\
\texttt{\$FF40} & LCDC & Write & Through \texttt{DISPLAY\_ON/OFF}, \texttt{SHOW/HIDE\_SPRITES}, \texttt{SHOW/HIDE\_BKG}, and \texttt{SHOW/HIDE\_WIN}; uses LCD enable bit 7, window enable bit 5, sprite enable bit 1, and background/window enable bit 0. \\
\texttt{\$FF47} & BGP & Write & Fade palettes \texttt{0xE4}, \texttt{0xF9}, \texttt{0xFE}, \texttt{0xFF}. \\
\texttt{\$FF4A} & WY & Write & \texttt{move\_win(MINWNDPOSX, DIALOGUE\_WIN\_Y)} positions the dialogue window at pixel Y 112. \\
\texttt{\$FF4B} & WX & Write & \texttt{move\_win} positions the dialogue window using \texttt{MINWNDPOSX}. \\
\bottomrule
\end{tabular}
\caption{Hardware registers touched by the source}
\end{table}

GBDK functions such as \texttt{set\_bkg\_data}, \texttt{set\_bkg\_tiles}, \texttt{set\_sprite\_data}, \texttt{set\_sprite\_tile}, and \texttt{move\_sprite} also write VRAM or OAM hardware areas. The project does not write the raw VRAM or OAM addresses manually.

\section{Interrupt Handlers}

\begin{description}
  \item[VBlank] Registered in \texttt{core} with \texttt{add\_VBL(vbl\_interrupt)}. Interrupts are enabled with \texttt{enable\_interrupts()}. The handler body is empty and contains comments stating it is called each frame when the PPU is in VBlank and should remain short.
  \item[Other interrupts] No LCD STAT, Timer, Serial, or Joypad ISR is registered by project code.
  % TODO: not found in codebase: non-VBlank interrupt handlers.
  \item[Cycle cost] No cycle count is documented for \texttt{vbl\_interrupt}.
  % TODO: not found in codebase: VBlank ISR cycle cost.
\end{description}
